\begin{helper}
Fix a two-bite embedding cell \(\pi\), red and blue center sets
\(B_R,B_B\), and red and blue projection sets \(P_R,P_B\).  Let
\[
K=\noderef{TwoBiteNaturalIndependenceScale}(1+\varepsilon,n),\qquad
p=\noderef{TwoBiteEdgeProbability}(1/2,n),
\]
and
\[
A=\frac{K n^{1/4-\varepsilon}}{6(\log n)^2},\qquad t=\lceil A\rceil.
\]
Form the normalized coordinate set from the red pairs
\(B_R\times P_R\) and blue pairs \(B_B\times P_B\), after discarding diagonal
pairs.  Suppose a candidate event \(C(\omega)\), inside the embedding cell
\(\noderef{TwoBiteEmbedding}(\omega)=\pi\), implies that at least \(t\) of
these normalized coordinates are present.  If the natural fixed-pair tail
inputs hold,
\[
0<A,\quad |P_R|\le K,\quad |P_B|\le K,
\]
\[
p\,((|B_R|+|B_B|)K)\le A/(4e),
\qquad
n^{3/4-2\varepsilon}\le A/4,
\]
then
\[
\noderef{TwoBiteEventProbability}
  (n,m,p,\{\omega:\noderef{TwoBiteEmbedding}(\omega)=\pi
     \text{ and } C(\omega)\})
\le
e^{-n^{3/4-2\varepsilon}}\,
\noderef{TwoBiteEventProbability}
  (n,m,p,\{\omega:\noderef{TwoBiteEmbedding}(\omega)=\pi\}).
\]
\end{helper}
\begin{proof}
Let \(D\) be the event that the embedding is \(\pi\), and let \(M\) be the
event that at least \(t=\lceil A\rceil\) of the normalized red/blue coordinate
set are present.  By the assumed implication, \(D\cap C\subseteq D\cap M\).
Since sample weights are nonnegative for \(0\le p\le1\), monotonicity of the
finite probability sum gives
\[
\mathbb P(D\cap C)\le \mathbb P(D\cap M).
\]
The embedding-cell coordinate union bound
\(\noderef{MediumClosedPairsEmbeddingCellCoordinateUnionBridge}\) gives
\[
\mathbb P(D\cap M)
\le
\binom{|B_R||P_R|+|B_B||P_B|}{t}p^t\,\mathbb P(D).
\]
The hypotheses listed above are exactly the inputs required by
\(\noderef{MediumClosedPairsNaturalFixedPairTail}\), so the binomial factor is
at most \(e^{-n^{3/4-2\varepsilon}}\).  Multiplying by the nonnegative cell
probability \(\mathbb P(D)\), supplied by
\(\noderef{TwoBiteEventProbabilityNonnegative}\), and chaining the two
inequalities proves the desired cell-tail estimate.
\end{proof}
